% ============================================================
%  Table 3 — Component Ablation (SVD-LLMv2 Table-4 style)
%  LLaMA-7B, WikiText-2, 60% compression ratio.
%  FR-SVD (F): Fisher-guided CEALC only
%  FR-SVD (R): Eigenspace residual compensation only
%  FR-SVD    : Both components
%  Required preamble packages:
%    \usepackage{booktabs}
%    \usepackage[table]{xcolor}
%    \usepackage{graphicx}   % for \resizebox
% ============================================================

\providecolor{mygreen}{rgb}{0.00,0.55,0.00}

\begin{table}[t]
\centering
\footnotesize
\renewcommand{\arraystretch}{1.15}

\caption{Perplexity ($\downarrow$) of LLaMA-7B compressed by SAES-SVD and
FR-SVD with individual\,/\,both components under 60\% compression ratio on
WikiText-2. FR-SVD (F) uses Fisher-guided CEALC truncation only; FR-SVD (R)
uses eigenspace residual compensation only; FR-SVD uses both.
Relative perplexity reduction against SAES-SVD is marked in
\textcolor{mygreen}{green}.}
\label{tab:ablation-components}

\resizebox{\linewidth}{!}{%
\begin{tabular}{c|c|c|c}
\midrule
SAES-SVD & FR-SVD (F) & FR-SVD (R) & FR-SVD \\ \midrule
22.01
 & 20.42\,(\textcolor{mygreen}{$\downarrow$7.2\%})
 & 21.13\,(\textcolor{mygreen}{$\downarrow$4.0\%})
 & \textbf{19.36}\,(\textcolor{mygreen}{$\downarrow$12.0\%}) \\ \midrule
\end{tabular}%
}
\end{table}
